\section{Appendix: proofs and auxiliary lemmas}

This appendix collects the auxiliary objects used in the minimax and complete-block arguments. The first group of statements isolates the one-block Hilbert-space calculations: a polynomial perturbation program for least-favourable directions, a dual weight program for block-local unbiased scores, and the representer identity tying both programs to the complete-block energy.

The block-local Riesz formulation works in the finite design space induced by the product Bernoulli law in \cref{obj:ass:bernoulli-design}. For a complete \(d\)-block, the low-order polynomial subspace \(\mathcal P_d\) contains raw monomials up to the exposed order, and the contrast \(L_d\) evaluates the all-one-versus-all-zero difference on that block.

\begin{definitionv}{def:perturbation-program}[Low-order block space $\mathcal P_d$, contrast $L_d$, and program $\mathsf{PerturbProg}_{d,\beta,p}$]
For \(d\ge 1\), let \(\mathcal P_d\) be the low-order polynomial subspace on \(d\) Bernoulli coordinates,
\[
\mathcal P_d
:=
\left\{
h(z)=\sum_{\substack{S\subseteq \{1,\ldots,d\}\\ |S|\le \bar\beta_d}}
a_S\prod_{j\in S}z_j
\right\}.
\]
Define the all-one-versus-all-zero block contrast functional \(L_d:\mathcal P_d\to\mathbb R\) by
\[
L_d(h)
:=
h(1,\ldots,1)-h(0,\ldots,0).
\]
The perturbation program is
\[
\mathsf{PerturbProg}_{d,\beta,p}
:=
\min_{h\in\mathcal P_d}\mathbb E_Z[h(Z)^2]
\]
subject to
\[
L_d(h)=1.
\]
\end{definitionv}

\Cref{obj:def:perturbation-program} identifies the minimum-energy direction that moves the block contrast by one unit. This is the direction used by the least-favourable complete-block construction in \cref{obj:def:block-family}: controlling its \(L^2(P_Z)\) energy controls the affinity between nearby prior-predictive laws while preserving a fixed separation in total treatment effects.

The dual calculation is a minimum-norm problem over square-integrable block weights. Its moment equations impose exact unbiasedness against every nonempty raw monomial that can enter a low-order block outcome.

\begin{definitionv}{def:weight-program}[Weight program \(\mathsf{WeightProg}_{d,\beta,p}\)]
For \(d\ge 1\), define
\[
\mathsf{WeightProg}_{d,\beta,p}
:=
\min_{w\in L^2(P_Z)}\mathbb E_Z[w(Z)^2]
\]
subject to
\[
\mathbb E_Z[w(Z)]=0
\]
and, for every nonempty \(S\subseteq\{1,\ldots,d\}\) with \(|S|\le\bar\beta_d\),
\[
\mathbb E_Z\left[w(Z)\prod_{j\in S}Z_j\right]=1.
\]
\end{definitionv}

\Cref{obj:def:weight-program} is the block analogue of the SNIPE unbiasedness equations. The feasible weights reproduce the all-treated-versus-all-control contrast on every exposed raw monomial, so the value of the program is the smallest block variance compatible with those unbiasedness restrictions.

The next lemma supplies the exact representer calculation behind both programs. It also records the energy scale and a uniform raw-coefficient-mass bound for the normalized representer coefficient \(h_{d,T}\), which is used in the block perturbations.

\begin{lemmav}{lem:block-energy-representer}[Block representer energy]
Fix an integer \(\beta\ge 1\) and a Bernoulli treatment probability \(p\in(0,1)\). Let \(k_\star(d,\beta,p)\) be the exposed order in \cref{obj:def:exposed-order}. For each \(d\ge 1\), write
\[
A_d
=
\sum_{r=1}^{\bar\beta_d}
\binom d r
\frac{\Delta_r(p)^2}{\{p(1-p)\}^r}.
\]
Let \(g_d\) denote the complete-block centered contrast score, \(h_d:=g_d/A_d\), and let \(h_{d,T}\) be the coefficient of the raw monomial indexed by \(T\subseteq\{1,\ldots,d\}\) in \(h_d\).

Then there exist constants \(c_1,c_2,H\in\mathbb R\), depending only on \(\beta\) and \(p\), such that \(0<c_1\le c_2\), \(0<H\), and for every \(d\ge 1\):
\begin{itemize}
\item \textbf{(Energy scale.)}
\[
c_1\binom d {k_\star(d,\beta,p)}
\le
A_d
\le
c_2\binom d {k_\star(d,\beta,p)}.
\]
\item \textbf{(Product Bernoulli identities.)} For every finite design \(D\) on \(\{0,1\}^d\) that is product Bernoulli with common probability \(p\) as in \cref{obj:ass:bernoulli-design}, and every \(f\in\mathcal P_d\),
\[
\mathbb E_D[g_d(Z)f(Z)]
=
L_d(f).
\]
Moreover,
\[
L_d(h_d)=1,
\qquad
\mathbb E_D[h_d(Z)^2]=A_d^{-1}.
\]
\item \textbf{(Perturbation program.)} The representer \(h_d\) is feasible for the perturbation program in \cref{obj:def:perturbation-program}. For every feasible perturbation \(h\),
\[
A_d^{-1}\le \mathbb E_D[h(Z)^2],
\]
with equality if and only if \(h=h_d\). Consequently, for any proof witnesses of \(0\le p\) and \(p\le 1\),
\[
\mathsf{PerturbProg}_{d,\beta,p}=A_d^{-1}.
\]
\item \textbf{(Weight program.)} For every weight \(w\) feasible for the weight program in \cref{obj:def:weight-program} under \(D\),
\[
A_d\le \mathbb E_D[w(Z)^2],
\]
with equality if and only if \(w=g_d\). Consequently, for any proof witnesses of \(0\le p\) and \(p\le 1\),
\[
\mathsf{WeightProg}_{d,\beta,p}=A_d.
\]
\item \textbf{(Coefficient mass.)}
\[
\sum_{T\subseteq\{1,\ldots,d\}} |h_{d,T}|
\le
H.
\]
\end{itemize}
\end{lemmav}

\Cref{obj:lem:block-energy-representer} is the algebraic core of the appendix. It establishes that the normalized complete-block score is the unique minimum-energy perturbation with unit block contrast, while the unnormalized score is the unique minimum-energy unbiased weight. The same calculation gives \(A_d\asymp_{\beta,p}\binom d{k_\star(d,\beta,p)}\), which is the local energy comparison used in \cref{obj:thm:degree-frontier,obj:thm:bounded-outcome-degree-frontier}.

The global SNIPE variance calculation needs one graph-combinatorial ingredient. The covariance terms are organized by shared assignment subsets across neighborhoods, and the next lemma converts those shared subsets into a single out-degree charge.

\begin{lemmav}{lem:overlap-count}[Overlap count bound]
Let \(V_n\) be the finite population and unit index set, let \(G\) be a directed interference graph on \(V_n\), and write
\[
N_i=\{j\in V_n:G(j,i)\}.
\]
Suppose:
\begin{itemize}
\item \textbf{(Bounded degree.)} The graph \(G\) satisfies the bounded-degree condition with bound \(d\) in \cref{obj:ass:bounded-degree}.
\item \textbf{(Order.)} The integer \(r\) satisfies \(r\ge 1\).
\item \textbf{(Unit.)} The unit \(i\) belongs to \(V_n\).
\end{itemize}
Then
\[
\sum_{l\in V_n}\binom{|N_i\cap N_l|}{r}
=
\sum_{\substack{S\subseteq N_i\\ |S|=r}}
\#\{l\in V_n:S\subseteq N_l\}
\]
and
\[
\sum_{\substack{S\subseteq N_i\\ |S|=r}}
\#\{l\in V_n:S\subseteq N_l\}
\le
d\binom{|N_i|}{r}
\le
d\binom{d}{r}.
\]
\end{lemmav}

\Cref{obj:lem:overlap-count} explains where the extra factor \(d\) enters the upper bounds. The binomial term \(\binom{|N_i|}{r}\) is already represented in the local block energy, while the count of units whose neighborhoods contain the same \(r\)-set is controlled by the out-degree part of \cref{obj:ass:bounded-degree}. This is the covariance-counting step used to pass from block energy to the global \(dA_d/n\) scale.

The remaining auxiliary statements support the lower bound. They place the two sign-indexed complete-block priors from \cref{obj:def:block-family} on a common dominating measure and bound the unhalved squared Hellinger distance between their densities.

\begin{lemmav}{lem:block-prior-dominating-measure}[Block prior integrates to one]
Use the notation for the Bernoulli assignment law and finite population from \cref{obj:ass:bernoulli-design}, the block size from \cref{obj:ass:bounded-degree}, the interaction order from \cref{obj:ass:low-order}, and the radius from \cref{obj:ass:bounded-coefficient-mass}. Let \(n,d,\beta\in\mathbb N\) and \(B,p,\sigma\in\mathbb R\), and set \(m=\lfloor n/d\rfloor\). Suppose that:
\begin{itemize}
\item \textbf{(Population and blocks.)} \(1\le n\) and \(1\le d\le n\).
\item \textbf{(Radius.)} \(B>0\), and let \(s=B/2\).
\item \textbf{(Design probability.)} \(p\in(0,1)\).
\item \textbf{(Prior sign.)} \(\sigma\in\{-1,+1\}\).
\end{itemize}
Let \(f_s\) be the cosine-squared baseline density from \cref{obj:def:block-family}, let \(\delta\) be the perturbation amplitude from \cref{obj:def:block-family}, and let \(h_d\) be the normalized complete-block representer from \cref{obj:def:block-score-energy}. Define the block dominating measure by
\[
\mu_{n,d}
=
\operatorname{count}_{\{0,1\}^{V_n}}\otimes \lambda^m .
\]
For \((z,u)\in\{0,1\}^{V_n}\times\mathbb R^m\), define the sign-\(\sigma\) block prior density by
\[
\pi_\sigma(z,u)
=
\Bigl(\prod_{j\in V_n} p^{z_j}(1-p)^{1-z_j}\Bigr)
\prod_{b=1}^{m}
f_s\bigl(u_b-\sigma\delta\,h_d(z_{V_b})\bigr),
\]
where \(V_1,\ldots,V_m\) are the complete \(d\)-blocks of \cref{obj:def:block-family}. Then
\[
\int_{\{0,1\}^{V_n}\times\mathbb R^m}\pi_\sigma(z,u)\,d\mu_{n,d}(z,u)=1.
\]
Thus \(\pi_\sigma\) is a probability density with respect to the common block dominating measure \(\mu_{n,d}\).
\end{lemmav}

\Cref{obj:lem:block-prior-dominating-measure} makes the prior comparison a calculation between two ordinary densities. The product Bernoulli assignment component is the same under both signs, while the observed block coefficients are shifted by the score-aligned perturbation from \cref{obj:def:block-family}.

\begin{lemmav}{lem:block-prior-hellinger-bound}[Block prior Hellinger bound]
Let \(n,d,\beta\) be integers and let \(B,p\in\mathbb R\). Suppose that:
\begin{itemize}
\item \textbf{(Design probability.)} The Bernoulli design notation and treatment probability \(p\) are as in \cref{obj:ass:bernoulli-design}, with \(0<p<1\).
\item \textbf{(Block size.)} The degree parameter \(d\) is as in \cref{obj:ass:bounded-degree}, with \(1\le d\le n\) and \(1\le n\).
\item \textbf{(Interaction order.)} The interaction-order parameter \(\beta\) is as in \cref{obj:ass:low-order}, with \(\beta\ge1\).
\item \textbf{(Radius.)} The radius \(B\) is as in \cref{obj:ass:bounded-coefficient-mass}, with \(B>0\).
\item \textbf{(Block priors.)} Let \(\mu_{n,d}\), \(\pi_{+1}\), and \(\pi_{-1}\) be the block dominating measure and the two signed block prior densities from \cref{obj:lem:block-prior-dominating-measure}.
\item \textbf{(Amplitude and energy.)} Let
\[
m:=\left\lfloor\frac nd\right\rfloor,
\qquad
\delta
:=B\min\{(2H_{\beta,p})^{-1},(4\pi)^{-1}\}
  \min\left\{1,\sqrt{\frac{A_d}{m}}\right\},
\]
where \(m\) and \(\delta\) are the complete-block count and perturbation amplitude from \cref{obj:def:block-family}, and \(A_d\) is the complete-block score energy from \cref{obj:def:block-score-energy}.
\end{itemize}
Then the unhalved squared Hellinger distance
\[
H^2(\pi_{+1},\pi_{-1})
:=\int\bigl(\sqrt{\pi_{+1}}-\sqrt{\pi_{-1}}\bigr)^2\,d\mu_{n,d}
\]
satisfies
\[
H^2(\pi_{+1},\pi_{-1})
\le
\frac{4\pi^2m\delta^2}{B^2A_d}.
\]
\end{lemmav}

\Cref{obj:lem:block-prior-hellinger-bound} supplies the affinity control for the two-point prior comparison. Combined with the total-effect separation stated in \cref{obj:thm:degree-frontier,obj:thm:bounded-outcome-degree-frontier}, it yields the saturated lower-bound branch through the complete-block construction and the linear branch through the choice of perturbation amplitude.

Together, these auxiliary statements organize the theorem proofs. The upper bounds in \cref{obj:thm:degree-frontier,obj:thm:bounded-outcome-degree-frontier} use the product-Bernoulli orthogonality in \cref{obj:lem:block-energy-representer} to identify SNIPE unbiasedness and then apply \cref{obj:lem:overlap-count} to bound the covariance sum by \(dA_d/n\). The lower bounds use the complete-block family in \cref{obj:def:block-family}, the coefficient-mass control in \cref{obj:lem:block-energy-representer}, and the density and Hellinger calculations in \cref{obj:lem:block-prior-dominating-measure,obj:lem:block-prior-hellinger-bound}. The complete-block local linear theorem is the finite-dimensional dual counterpart: \cref{obj:def:weight-program} and \cref{obj:lem:block-energy-representer} identify the least possible block energy for unbiased weights, and \cref{obj:thm:sharp-local-linear-constant-and-representers} aggregates that block criterion across complete blocks.
